% Filled in from the actual runs (see evidence/).
\newcommand{\conlecheSummary}{我们还在本机用独立检查器 con-leche 重放了完整的导出证明：1,468,418 个声明全部被接受，五个主定理全部被判定为 ``justified by the solution''，
而用于比对的挑战导出是我们自己从官方 Mathlib 重建并逐字节核对过的。}
\newcommand{\conlecheDetail}{%
\begin{verdictbox}
\textbf{结果：通过。}退出码 0。标准输出全文如下：
\VerbatimInput[fontsize=\scriptsize,breaklines=true,breakanywhere=true]{../evidence/con-leche/compare.stdout.log}
\end{verdictbox}

\noindent 即：解答导出的 1,468,418 个声明全部在 \lean{--verified} 模式下被接受；挑战导出的 51,692 个声明被接受，
其中 5 个以 \lean{sorry} 占位的定理——包括 \lean{riemannHypothesis}——全部由解答「证成」。这两个数字与 MIT 页面公布的 con-leche 结果完全一致。
本机运行耗时 73 分 9 秒（解析 340.3 秒，检查 4022.3 秒，16 个工作线程），峰值内存 47.4\,GB；
单个声明中最慢的是 \lean{GlobalTheory.HeightCounting.SideCopy.LamCopy.lam_copy}（374 秒）。日志见 \path{evidence/con-leche/}。

这一步同时回答了两个问题：其一，证明项确实通过了一个独立于官方内核实现的检查器（因而 \S\ref{sec:hygiene} 中 \lean{kexactB}／\lean{kcopy} 所依赖的「内核把关」是经得起独立复核的）；
其二，被证明的五个命题与「只导入 Mathlib 的五个陈述」是同一组命题——连同它们读到的每一个定义。
}
\newcommand{\scanThreeSummary}{第三遍扫描（结构体字段读取与消去）的要点：\lean{OAt.rec}、\lean{DivT.rec}、\lean{DivX.rec}、\lean{DV.Geo.rec} 均为零引用，
即这些记录的内容只可能经字段投影读出；\lean{OAt} 的 \lean{geo} 字段没有转手以外的读取者；内容提取引理 \lean{GeoTiedP.geoTied}、\lean{GeoTiedP.formsTieX} 不在依赖锥内；
而抽象接口 \lean{DV.Geo} 的字段确被锥内 621 个声明读取（\S\ref{sec:fact-degenerate}）。}
\newcommand{\scanThreeAppendix}{%
\subsection{第三遍扫描：结构体字段的读取与消去}

「依赖数」是传递依赖该模式的声明个数；投影类种子匹配的是内核项中的 \lean{proj} 结点。

\begin{center}\small
\begin{tabular}{@{}lrc@{}}
\toprule
种子（模式） & 依赖它的声明数 & 主定理依赖它？\\
\midrule
\tableinput{tables/scan3_counts.tex}
\bottomrule
\end{tabular}
\end{center}

直接读取 \lean{OAt} 的 \lean{lam} 字段的定理有三个：\lean{divEx_geoCl}、\lean{oAt_setDiv}（二者是转手）以及 \lean{GlobalTheory.T.parameter_of_discrete_series_ray_alt}（真正的消费者）；
直接读取 \lean{geo} 字段的只有前两个。实际出现在依赖锥内的 \lean{DV.Geo} 字段访问函数共 18 个：
\lean{Var}、\lean{Cl}、\lean{Dv}、\lean{dvGrp}、\lean{Fn}、\lean{fnField}、\lean{fnAlg}、\lean{H0}、\lean{supp}、\lean{dim}、\lean{Adl}、\lean{ofW}、\lean{inter}、\lean{twist}、\lean{vanishOn}、\lean{arTop}、\lean{arVol}、\lean{nefA}。
}

